\centerline{\bf \Huge Suiiiiмедианы}
\centerline{\bf \Huge }

\problem{1}
Сторона $A C$ треугольника $A B C$ является хордой окружности, касающейся прямой $B C$ в точке $C$. Вторая окружность проходит через точки $B, C$ и касается прямой $A C$ в точке $C$. Пусть $D \neq C$ - точка пересечения этих окружностей. Докажите, что $C D$ - симедиана треугольника $A B C$.

\problem{2}
Через вершины $A$ и $B$ треугольника $A B C$ проведена окружность, касающаяся прямой $B C$, а через вершины $B$ и $C$ - другая окружность, касающаяся прямой $A B$. Продолжение общей хорды $B D$ этих окружностей пересекает отрезок $A C$ в точке $E$, а продолжение хорды $A D$ одной окружности пересекает другую окружность в точке $F$. Найти отношение $A E: E C$, если $A B=5$ и $B C=9$. Сравнить площади треугольника $A B C$ и $A B F$.

\problem{3} 
Треугольник $A B C$ вписан в окружность с центром $O$. Окружность, построенная на $B O$ как на диаметре, повторно пересекает описанную окружность треугольника $A O C$ в точке $S$. Докажите, что $B S$ - симедиана треугольника $A B C$.

\problem{4}
На сторонах $A B$ и $B C$ треугольника $A B C$ во вне построены квадраты $A B X Y$ и $B C K L$ Прямые $X Y$ и $K L$ пересекаются в точке $T$. Докажите, что $B T$ - симедиана треугольника $A B C$.

\problem{5} 
Даны окружность, её хорда $A B$ и середина $C_0$ меньшей дуги $A B$. На большей дуге $A B$ выбирается произвольная точка С. Касательная к окружности, проведённая из точки $C$, пересекает касательные, проведённые из точек $A$ и $B$, в точках $X$ и $Y$ соответственно. Прямые $C_0 X$ и $C_0 Y$ пересекают прямую $A B$ в точках $N$ и $M$ соответственно. Докажите, что длина отрезка $N M$ не зависит от выбора точки $C$.

\problem{6} 
Внутри равнобедренного треугольника $A B C$ выбрана точка $P$ такая, что $\angle P B C=$ $\angle P C A$. Пусть $M$ - середина основания $B C$. Докажите, что $\angle B P M+\angle C P A=180^{\circ}$.

\problem{7}
В треугольнике $A B C$ проведена биссектриса $B D$ (точка $D$ лежит на отрезке $A C$ ). Прямая $B D$ пересекает окружность $\Omega$, описанную около треугольника $A B C$, в точках $B$ и $E$. Окружность $\omega$, построенная на отрезке $D E$ как на диаметре, пересекает окружность $\Omega$ в точках $E$ и $F$. Докажите, что прямая, симметричная прямой $B F$ относительно прямой $B D$, содержит медиану треугольника $A B C$.